\chapterwithopening{Software and Engineering Implementation}{chap:implementation}{
  \chaptercontentsitem{sec:implementation_introduction}
  \chaptercontentsitem{sec:implementation_environment}
  \chaptercontentsitem{sec:implementation_system_architecture}
  \chaptercontentsitem{sec:web_application_implementation}
  \chaptercontentsitem{sec:backend_database_storage_implementation}
  \chaptercontentsitem{sec:authentication_authorization_security}
  \chaptercontentsitem{sec:role_based_multi_organization_access}
  \chaptercontentsitem{sec:feature_availability_extension_implementation}
  \chaptercontentsitem{sec:real_time_communication_live_sessions}
  \chaptercontentsitem{sec:ai_learning_support_ide_implementation}
  \chaptercontentsitem{sec:development_workflow_deployment}
  \chaptercontentsitem{sec:implementation_challenges_adopted_solutions}
  \chaptercontentsitem{sec:implementation_chapter_summary}
}

\section{Introduction}\label{sec:implementation_introduction}

This chapter documents the verified software and engineering implementation of
EduVerse based on the current repository, the implementation analysis report,
and the accompanying project notes. It explains how the platform was realized
as a web-first educational workspace and how its main services, security
controls, role model, and modular academic tools were integrated into one
coherent system.

The chapter is deliberately limited to what is supported by the available
implementation evidence. The current repository clearly contains the web
application and its supporting server logic. It does not contain the mobile
companion source code, although project documentation references that mobile
client as a separate sibling repository. Accordingly, the discussion below
describes the shared web-facing implementation in detail and refers to mobile
integration only where the repository explicitly supports it through shared API
behavior.

\section{Implementation Environment}\label{sec:implementation_environment}

EduVerse was implemented as a full-stack web application rather than as a thin
frontend connected to a separately maintained backend service. The browser
client, server-side route handlers, shared domain logic, database access, and
external service integrations are coordinated within the same engineering
environment.

\begin{table}[htbp]
  \centering
  \small
  \renewcommand{\arraystretch}{1.15}
  \setlength{\tabcolsep}{5pt}
  \caption{Main EduVerse Implementation Environment}
  \label{tab:eduverse_implementation_environment}
  \begin{tabular}{|p{0.24\textwidth}|p{0.27\textwidth}|p{0.39\textwidth}|}
    \hline
    \textbf{Implementation layer} & \textbf{Main technologies} & \textbf{Implementation purpose} \\
    \hline
    Web client & Next.js App Router, React, TypeScript & Role-specific dashboards, class workspaces, forms, navigation, interactive learning tools, and client-side state. \\
    \hline
    Server/API layer & Next.js route handlers, shared domain logic in \texttt{lib/*}, validation helpers & Authenticated reads, mutations, workflow rules, file handling, integration calls, and reusable server-owned APIs. \\
    \hline
    Data and authentication & Supabase Postgres, Supabase Auth, RPC functions, row-level security & Persistent academic records, identity management, role selection, guarded data access, and workflow-level database operations. \\
    \hline
    File storage & AWS S3 & Storage of large educational files such as materials and assignment files, with signed retrieval. \\
    \hline
    Real-time and live sessions & LiveKit, Supabase Realtime & Live session entry, session presence coordination, and live update synchronization. \\
    \hline
    AI support & OpenRouter, PDF/image extraction helpers, browser-side AI interfaces & Class AI assistance, study summaries, drafting support, and contextual academic help. \\
    \hline
    Messaging and invite delivery & Gmail API when configured & Organization invite delivery through email when the required server configuration is available. \\
    \hline
    Review and deployment support & Bun test execution, GitHub workflow, GitHub Actions, Vercel preview deployments & Automated checks, branch review, and safe pre-merge inspection of feature increments. \\
    \hline
  \end{tabular}
\end{table}

\FloatBarrier

The implementation environment also reflects an intentional client asymmetry.
The web application is the primary and fully evidenced interface in the current
repository, while the mobile companion is treated as a separate client expected
to consume the same protected API layer rather than re-implement core business
logic independently.

\section{System Architecture}\label{sec:implementation_system_architecture}

The implemented architecture follows a web-first, service-integrated model. The
browser renders role-sensitive application pages, while most guarded reads and
mutations pass through server-side route handlers under \texttt{app/api/*}.
Those routes coordinate with Supabase for authentication, relational data,
database RPC functions, and some realtime-aware state. Specialized services are
then used only where they add a distinct capability that is not naturally
stored inside the relational database.

At a high level, the implemented request flow can be described as follows:

\begin{itemize}
  \item The browser loads the Next.js application and renders role-specific
    dashboard and class pages.
  \item The application shell resolves the authenticated user and the active
    organization/role context before protected workspaces are opened.
  \item Client actions call server-owned API routes, which apply validation,
    authorization checks, and workflow rules.
  \item The API layer reads and writes structured academic data through
    Supabase tables and, where appropriate, invokes Supabase RPC functions for
    higher-level workflow control.
  \item Large files, live-session access, and AI generation are delegated from
    the server layer to AWS S3, LiveKit, and OpenRouter respectively.
\end{itemize}

This architecture solves an important consistency problem. Because the platform
centralizes business rules in shared server-owned routes and helpers, the same
authorization rules and academic data flows can be reused by the web
application, by the separate mobile companion, and by future controlled
integrations without duplicating the role model in multiple clients.

The repository structure also reflects modular engineering. Application routes,
feature modules, shared libraries, and service helpers are separated so that
feature slices such as assignments, exams, live sessions, AI, and IDE behavior
can evolve while still depending on a shared organization-aware core.

\section{Web Application Implementation}\label{sec:web_application_implementation}

The EduVerse web application is the main operational client of the system. It
implements the complete educational workspace visible in the current
repository, including organization administration, class participation,
materials, communication, assignments, exams, results, live sessions, AI
support, and the browser-based coding lab.

The interface is organized around role-specific dashboards and class-centered
workspaces. Organization Administrators receive organization-scoped views for
class oversight, membership management, public access, feature configuration,
teacher permissions, archived terms, and settings. Teachers receive dashboards
and class tools oriented toward teaching activity, grading workload, and class
management. Students receive dashboards and class views focused on active
learning tasks such as materials, submissions, results, notifications, live
session participation, and available study tools.

At the class level, the web application acts as a unified academic workspace
rather than a simple content portal. Chat, announcements, uploaded materials,
assignment workflows, exam workflows, results visibility, live sessions,
whiteboard interaction, AI support, and IDE access are all exposed within the
same class environment when enabled. This reduces the fragmentation that would
otherwise occur if each activity were moved to a different external tool.

The same class structure is also flexible in purpose. A class can operate as a
conventional taught course, but the implementation also supports narrower uses,
such as a library-style class dedicated mainly to sharing books and learning
resources. This confirms that EduVerse is not limited to a single pedagogical
format; it provides a governed academic workspace that can be specialized
according to organizational needs.

Although the broader project includes a mobile companion direction, the current
repository proves the web application only. For that reason, this chapter treats
the web client as the verified primary implementation and does not claim proven
mobile feature parity from the available source tree.

\section{Backend, Database, and Storage Implementation}
\label{sec:backend_database_storage_implementation}

EduVerse uses a Supabase-centered backend model implemented from within the
Next.js project. No separate standalone backend service was identified in the
repository. Instead, the route handlers and shared server libraries perform the
main backend responsibilities: access control, request validation, workflow
execution, data persistence, and integration with external services.

The relational data model is organized into several main domains: organization
membership and access, class management, feature and extension settings,
materials and messages, assignments, exams, notifications, live-session
records, and historical academic data. Route handlers query these tables
directly where appropriate and also invoke Supabase RPC functions for higher
level actions such as organization creation, selected-role updates, class
management checks, join-link acceptance, results visibility updates, and
live-session claiming.

An important design choice concerns file handling. Large media and document
objects are not stored directly inside the frequently queried relational tables.
Instead, EduVerse keeps structured academic metadata in the database while
storing uploaded files in AWS S3. Materials, assignment attachments, and
student-submitted files therefore preserve their class ownership, uploader,
access, and descriptive metadata in Supabase, while the binary content itself
is stored separately in object storage. This separation reduces unnecessary
load on the transactional database and keeps routine academic queries focused
on structured records rather than large file payloads.

File handling is also server-mediated. The browser sends uploads to the server,
the server validates the request, writes the file to S3 using protected
credentials, and then stores the related metadata in the database. Download
access is issued through short-lived signed URLs, and some material content is
also proxied through authenticated server routes for inline viewing. This
approach keeps storage credentials outside the browser and aligns file delivery
with the same access-control model used for the rest of the platform.

Selective privileged access is used where necessary. The repository includes a
service-role Supabase client for restricted server-only operations such as
audit-log writes and selected aggregation or registration tasks. Even so, the
general pattern remains consistent: privileged access is limited, while most
application behavior stays inside request-scoped authenticated route handling.

\section{Authentication, Authorization, and Security}
\label{sec:authentication_authorization_security}

Authentication in EduVerse is built on Supabase Auth. The repository shows
sign-in, sign-up, password reset, password change, authentication callbacks,
and protected route entry behavior. Once the user is authenticated, the system
still requires organization and role context before organization-scoped or
class-scoped pages may be opened.

Authorization is implemented in multiple layers rather than through one
front-end check alone. The application shell prevents unauthenticated access to
protected areas, route handlers enforce access through server-side user
resolution, selected-role records determine which role the user is currently
operating under, helper logic and RPC functions perform class-level permission
checks, and database migrations define row-level security policies across major
tables. This layered model is important for a platform in which the same person
may hold different roles in different organizations.

The implementation also contains several concrete security controls:

\begin{itemize}
  \item server-side validation for uploads, assignments, exams, and other
    guarded inputs before persistence;
  \item short-lived signed S3 URLs for protected file downloads;
  \item invite acceptance tied to the authenticated email address rather than
    token possession alone;
  \item public join links with enable/disable state, expiry, usage limits,
    approval options, and role targeting;
  \item privacy filtering for student-facing results summaries;
  \item exam passcode hashing, invalid-attempt cooldown handling, integrity
    event recording, flagging, voiding, and audit logging.
\end{itemize}

The repository did not show webcam-based proctoring, biometric checks, or
direct browser-owned upload credentials. Security is therefore centered on
authenticated access, row-level policies, route protection, signed file
delivery, and event-based integrity monitoring rather than on intrusive
surveillance mechanisms.

\section{Role-Based and Multi-Organization Access}
\label{sec:role_based_multi_organization_access}

One of the central implementation problems solved by EduVerse is the need for a
single user to belong to multiple organizations and to operate under different
roles inside one organization. The repository addresses this through
organization-aware membership records, multi-role membership support, active
role selection, and role-filtered API behavior.

In the underlying implementation, the Organization Administrator role uses the
legacy internal identifier \texttt{org\_admin}. In the thesis wording,
however, it is more accurate to describe this role as the Organization
Administrator because that title better reflects its institutional
responsibilities. Those responsibilities include organization setup, user and
invite management, public access configuration, teacher permissions, class
oversight, archived term review, and organization-level settings.

Multi-role support is realized through dedicated membership-role records and a
selected-role pointer. A user can therefore hold more than one role within the
same organization, but the active interface and permitted operations are always
resolved through the currently selected role. This prevents the platform from
collapsing all privileges into a single undifferentiated identity and makes
role transitions explicit and auditable.

The same access model continues at class level. Teacher activity is constrained
by organization permissions and class assignment. Student access depends on
class membership, active role, feature settings, and visibility rules.
Organization-visible classes and controlled public join workflows expand access
only within the boundaries defined by organization policy. Archived classes and
previous-term data are likewise preserved without abandoning the same access
structure.

\section{Feature Availability and Extension Implementation}
\label{sec:feature_availability_extension_implementation}

EduVerse was designed as a configurable platform rather than as a fixed bundle
in which every organization and every class must use the same tool set. The
implementation report shows this clearly through feature definitions, presets,
organization-level feature settings, class-level feature settings, organization
extensions, and class extension settings.

This configuration model serves two purposes. First, it reduces interface and
workflow clutter by hiding or blocking tools that are not enabled for the
current organization or class. Second, it allows the same academic platform to
serve different institutional policies and teaching styles without fragmenting
into separate products. The AI feature is a strong example of this approach:
AI-assisted help, summaries, drafting, and feedback support can be enabled
where useful and disabled where instructors or organizations consider it
counterproductive.

The same flexibility extends to extensions and special-purpose classes. Because
tool exposure is governed through organization and class settings, a class can
be configured around a narrow purpose, such as a resource-sharing or
library-oriented workspace, without forcing unrelated assessment or coding
tools into the user experience. Similarly, public organizations and public join
workflows can support broader sharing and open-learning access while still
preserving role constraints, approval logic, and organization-level control.

From an engineering perspective, this modular feature model also improves
integration potential. Since core behavior is exposed through shared protected
API routes rather than through web-only client assumptions, future clients or
external institutional integrations can reuse the same academic data model and
authorization rules without duplicating the platform's business logic.

\section{Real-Time Communication and Live Sessions}
\label{sec:real_time_communication_live_sessions}

EduVerse implements both asynchronous and synchronous communication inside the
same class workspace. Asynchronous interaction is provided through class chat,
announcements, notifications, and message-linked media. Synchronous interaction
is provided through live-session entry, session presence control, and a shared
whiteboard environment.

The live-session implementation integrates LiveKit and Supabase Realtime. The
server issues LiveKit access tokens for authorized participants, while
Supabase-backed session records are used to claim the active session, update
heartbeats, recover from stale sessions, and synchronize state in the client
store. This avoids treating live teaching as an isolated external meeting and
keeps it connected to the same class membership and access rules used
elsewhere in the platform.

The whiteboard implementation further strengthens this integration. The
repository evidences synchronized whiteboard behavior with strokes, shapes,
text, viewport synchronization, clear, undo, and redo actions. At the same
time, the database evidence shows no dedicated whiteboard-state table, which
suggests that live collaboration behavior is coordinated primarily through the
session layer rather than through a separate long-term whiteboard persistence
schema.

\section{AI Learning Support and IDE Implementation}
\label{sec:ai_learning_support_ide_implementation}

AI support in EduVerse is implemented as a contextual academic service rather
than as an unrestricted external chatbot. The repository shows class AI agent
routes, material extraction and summary generation, assignment-related support,
exam drafting assistance, and AI-backed feedback support. These flows gather
class context from the system before generating responses through OpenRouter.

The implementation is also role-aware. Student-facing AI behavior is bounded by
prompt policies intended to avoid direct final answers during active assessed
work, while teacher-facing AI support is more oriented toward drafting,
preparation, and assistance. Because AI is treated as one governed feature among
others, it can be disabled at organization level or class level when a learning
context requires a more restrictive policy.

The integrated coding lab follows the same embedded-tool philosophy. EduVerse
implements a browser-based programming workspace with Monaco editor support,
file-tree management, previews, templates, a virtual terminal, and browser-safe
runners. The repository does not show real operating-system process execution
for the IDE. Instead, the implementation deliberately avoids unsafe direct
command execution and keeps the coding lab inside a browser-safe execution
model. This makes the current IDE appropriate for guided educational practice,
while more powerful execution remains a future enhancement rather than an
unsupported claim.

\section{Development Workflow and Deployment}
\label{sec:development_workflow_deployment}

The implementation approach described in Chapter~3 was supported by a practical
engineering workflow that kept feature increments reviewable and safe to
integrate. Project development notes indicate a GitHub-based fork, branch, and
pull-request workflow in which individual changes were developed in bounded
feature branches and reviewed before merge. This process aligned well with the
project's vertical-slice methodology because each increment could be evaluated
as one coherent academic feature rather than as disconnected interface and
backend fragments.

Automated checks also supported this workflow. GitHub Actions was used to run
build, formatting, and type-consistency checks before integration, while Vercel
preview deployments provided a practical environment for reviewing UI behavior
and cross-page interactions before a change entered the main branch. These
deployment previews were especially helpful for a platform such as EduVerse,
where role-sensitive behavior and feature-gated pages often need visual review
in addition to code inspection.

The shared server-owned API layer further contributed to workflow stability.
Because the same protected routes could be reused across the web client and the
separate mobile companion direction, the project avoided scattering business
rules across multiple client implementations. That engineering decision reduces
future maintenance effort and makes later client expansion more manageable.

\section{Implementation Challenges and Adopted Solutions}
\label{sec:implementation_challenges_adopted_solutions}

The implementation of EduVerse required more than the accumulation of isolated
features. Several recurring engineering challenges had to be addressed so that
the platform would remain coherent across roles, organizations, storage needs,
and review workflows.

\begin{table}[htbp]
  \centering
  \footnotesize
  \renewcommand{\arraystretch}{1.15}
  \setlength{\tabcolsep}{4pt}
  \caption{Implementation Challenges and Adopted Solutions in EduVerse}
  \label{tab:implementation_challenges_solutions}
  \begin{tabular}{|p{0.28\textwidth}|p{0.60\textwidth}|}
    \hline
    \textbf{Implementation challenge} & \textbf{Adopted solution} \\
    \hline
    One user needing access to multiple organizations & Active organization selection, organization-aware API behavior, and organization-scoped payload loading through the authenticated user context. \\
    \hline
    One membership needing more than one role & Separate membership-role records, selected-role storage, role-switching UI, and route checks tied to the current active role. \\
    \hline
    Keeping web and mobile behavior consistent & Reuse of server-owned authenticated API routes and bearer-token-aware request handling instead of duplicating business rules per client. \\
    \hline
    Storing large educational files without weakening relational queries & Separation of binary file storage into AWS S3 while keeping structured metadata, ownership, and access rules in Supabase tables. \\
    \hline
    Supporting public access without losing governance & Public join links, approval-required modes, expiration controls, role targeting, visibility rules, and organization-level public-feature settings. \\
    \hline
    Handling live-session concurrency and reconnect behavior & Session claiming, heartbeat updates, stale-session recovery windows, synchronized live-session records, and realtime updates. \\
    \hline
    Preserving exam fairness and traceability & Passcode protection, exam lock mode, integrity-event capture, audit logging, flagged/voided attempt states, retakes, and explicit results release. \\
    \hline
    Allowing different organizations and classes to use different tool sets & Feature registry, organization/class feature settings, presets, extension settings, and controlled hiding/blocking of unavailable tools. \\
    \hline
    Integrating team contributions safely & GitHub branch and pull-request review, automated GitHub Actions checks, and Vercel preview deployments before merge. \\
    \hline
  \end{tabular}
\end{table}

\FloatBarrier

These adopted solutions show that the platform's main technical contribution is
not only the presence of features, but the way those features were organized so
that governance, access control, academic data, live interaction, storage, and
review workflow all remained compatible with one another.

\section{Chapter Summary}\label{sec:implementation_chapter_summary}

This chapter explained how EduVerse was implemented as a web-first,
multi-organization educational workspace built around shared server-owned APIs,
Supabase-backed data and access control, S3-based file storage, LiveKit-backed
live sessions, OpenRouter-based AI support, and a browser-safe integrated
coding lab. It also showed how modular feature control, role-aware security,
and disciplined review/deployment practices were used to keep the platform
coherent as it expanded. Chapter~7 evaluates this implementation through the
available testing and validation evidence.
